Completar um look é sugerir uma peça de outra categoria que combine com as peças já escolhidas. A literatura resolve esse problema com redes neurais de compatibilidade, como a Bi-LSTM e os espaços de embeddings por tipo de peça, e com a recomendação top-$N$ baseada em itens. Este trabalho mede até onde estruturas de dados clássicas, que podem ser explicadas passo a passo, resolvem a mesma tarefa. A base é a FashionVerse-MME, com 102.370 peças e 31.333 looks. Uma KD-tree acha as peças parecidas. Uma BFS de dois passos no grafo de peças e looks conta os looks que duas peças têm em comum. Os grupos de estilo saem do algoritmo de Kruskal com Union-Find, e o algoritmo de Tarjan aponta as articulações do grafo entre grupos. A avaliação dá nota à categoria inteira, com 10 ordens de desempate e intervalos de bootstrap. No teste, com 5.842 looks, a versão final F teve Hit@10 de 1,056\% e FITB de 32,04\%, contra 0,182\% e 25,00\% do sorteio e 0,751\% e 25,37\% da versão C. A diferença para a C é real nas duas métricas. Contra a popularidade, a vantagem no Hit@10 pode ser sorte. A principal limitação é a peça nova: em 72,6\% dos looks de teste a peça escondida não aparece no treino, e nesse grupo a F não acerta nenhuma no top 10. Com a peça conhecida, a F acerta 3,856\% no Hit@10, 23,0 vezes o sorteio.
